\section{Legal Domain}
\label{sec:domains}
In this section, we describe the Italian legal system, the different types of trials with the documents involved, and the relationship between Italian law, \acl{ai} and predictability (ai act, gdpr).
italian criminal proceedings, vs civil
\todo{introdurre ai prima}
\todo{different law codes, civil law, supreme court (vs common law)}

The Italian Republic uses the \emph{civil law} system, in contrast with countries such as the United Kingdom of Great Britain and Northern Ireland and the United State of America, using the \emph{common law} system. Civil law is based on codes, while common law is based on case-law~\cite{Carlson2009}. Typically, civil law systems rely on four foundational codes: civil code, the civil procedure code, criminal code, and criminal procedure code. \todo{here add a brief descriptions of what are these codes} These codes are intended as coherent and consistent principles for supporting the broader legal system. In contrast, common law systems typically develop codes in response to specific issues, often drafting them only after principles have already emerged through case law~\cite{Carlson2009}.
Furthermore, in common law, precedent cases are of central importance, as judges are required to follow earlier decisions, and a single ruling from a higher court may determine the outcome of a case. This is different in civil law, where a single precedent case is generally not treated as binding, with the possible exception of decisions in the constitutional sphere~\cite{Carlson2009}. 

The organizational structure of civil law court systems is broadly similar to that of common law systems, featuring a court of first instance, a court of appeals, and a supreme court. Additionally, specialized courts exist with jurisdiction limited to particular areas of law. Unlike in common law jurisdictions, juries are generally not used in civil law systems. The cases are decided by judges, eventually including lay judges~\cite{Carlson2009}, such as in the Italian court of assize~\cite[Article 3]{legge10apr1951}---which is competent to judge serious criminal offenses, e.g., the ones carrying maximum penalties of at least 24 years of imprisonment~\cite[Article 5]{CPP}.  

\todo{cassazione lavora sulle leggi non sul fatto. parlare di come funzionano i gradi citando articoli}



\section{Inception of Legal Proceedings in the Italian Legal System}

In the Italian legal system, a legal proceeding (\textit{procedimento}) may arise in either the civil or the criminal sphere, with distinct initiation mechanisms reflecting the nature of the dispute.

\subsection{Civil Proceedings}

Civil proceedings (\textit{procedimenti civili}) are \emph{initiated by a private party}. The \emph{plaintiff} begins the process by filing an \textit{atto di citazione} (writ of summons) with the competent court. This document must indicate: the identities of the parties, the tribunal being addressed, the factual and legal basis of the claim, the evidence intended to be produced, and the date of the hearing \cite[Article 163]{CPC}. 

The \emph{defendant} responds with a \textit{comparsa di risposta} (statement of defense), outlining their defenses and any counterclaims \cite[Article 167]{CPC}. From this moment, the judicial authority schedules hearings, manages the collection of evidence, and ultimately renders a decision (\textit{sentenza}) which must include motivations as required by the Constitution~\cite[Article 111]{Costituzione}.

\subsection{Criminal Proceedings}

Criminal proceedings (\textit{procedimenti penali}) begin with the occurrence of a potential criminal offense. When a \emph{notitia criminis} (report of a crime) is received---via complaint, denunciation, or ex officio observation\todo{meaning}---the \textit{Pubblico Ministero} opens the preliminary investigation (\textit{indagini preliminari})~\cite[Article 326]{CPP} under the authority of the \textcite{CPP}. \todo{aggiungi citazioni ai codici direttamente. poi agli articoli}

During this phase, investigative acts are performed by the \textit{Pubblico Ministero} and the \textit{Polizia Giudiziaria}, and documented in official records (\textit{verbali}) as required by \cite[Article 373]{CPP}. The suspect is formally notified through an \textit{informazione di garanzia} \cite[Article 369]{CPP}. If sufficient evidence exists, the prosecutor requests an indictment (\textit{richiesta di rinvio a giudizio}) and the judge for the preliminary hearing (\textit{GUP}) decides whether the case proceeds to trial \cite[Article 416]{CPP}.




\todo{something general for legal}

\cite{BELLANDI2024,pozzi_wise_2025,aixia_2023,Agazzi2025DAVE,real_batini2021semantic}
\todo{italian legal system. civile penale, investigations. usa tesi dottorato legal}
\todo{AI act. news for legal AI italy, batini,...}
\subsection{Criminal Investigations}
Data in criminal investigations comes from a heterogeneity of sources such as documents, e.g., reports by consultants, (semi-)structured files, e.g., phone call records, bank transfers, and digitalized papers, e.g., seized in the investigation.
\todo{add example}

\todo{focus on IMA}

\begin{figure}[htbp]
\centering
\begin{minipage}{0.92\linewidth}
\ttfamily
\small
\textbf{Chat Transcript (Extract)}\\
Participants: Mario Rossi (M), Luca Bianchi (L)\\
Date Range: 2025-07-28 to 2025-07-30\\[0.5em]
\textbf{[2025-07-28 -- 14:12]} M: Hey, are we still meeting on Friday?\\
\textbf{[2025-07-28 -- 14:13]} L: Yeah, that's August 1st, right?\\
\textbf{[2025-07-28 -- 14:14]} M: Correct.\\
\textbf{[2025-07-28 -- 14:15]} L: Where though? The place near Piazza Vittorio Veneto in Turin or the one by the river?\\
\textbf{[2025-07-28 -- 14:17]} M: Let's do the bar at Via Roma 210, easier to park.\\
\textbf{[2025-07-28 -- 14:18]} L: Fine. Anyone else coming?\\
\textbf{[2025-07-28 -- 14:19]} M: Maybe Gianni, but he's in Milan until Thursday.\\[0.5em]
\textbf{[2025-07-29 -- 08:45]} M: \textit{[Voice message -- 2 min 14 sec]}\\
\quad \textit{Transcription:}\\
\quad So listen, I was thinking\dots it might be better if we meet a bit earlier, like around 6:30 PM, because I have to talk to you about something before Gianni gets there. It's about that thing we discussed last time at Marco's place — you know, the one near Porta Susa station. I also want to avoid too many people seeing us together at the main bar during the busy hour. If we meet earlier, we can grab a table at the back, order something small, and then, when Gianni arrives, we just act like we've been casually catching up. Also, there's a parking spot I found behind the cinema on Via Po, usually empty at that time. Oh, and bring the envelope I gave you last month, I'll need to check something inside. Don't text me too much about this, we'll talk in person.\\[0.5em]
\textbf{[2025-07-29 -- 09:02]} L: Got it. 6:30 PM, Via Roma 210, Friday. I'll bring the envelope.\\
\textbf{[2025-07-30 -- 16:47]} M: Perfect.
\end{minipage}
\caption{Simulated chat exchange containing dates, locations, persons, and a long voice message transcription.}
\label{fig:chat-investigation-example}
\end{figure}


\subsection{Judgments}
Legal documents, and especially court judgments and orders, contain information that is valuable in multiple applications (from legal case retrieval to legal process analysis), for different purposes (from the comparison with similar cases for uniform judgments to discovering trends for specific legal subjects, e.g., counting the average maintenance in relation to economical conditions of the partners in divorces), and stakeholders (from judges to law administrators, lawmakers, lawyers, scholars, and the general public).
\todo{add example in english}

\begin{figure}[htbp]
\centering
\begin{minipage}{0.9\linewidth}
\ttfamily
\small
ORDINARY COURT OF TURIN\\
Civil Division\\
Judgment No. 1234/2025\\
Published on March 15, 2025\\[1em]
ITALIAN REPUBLIC\\
IN THE NAME OF THE ITALIAN PEOPLE\\[1em]
The Court, sitting as a single judge, in the person of Judge Andrea Bianchi, has delivered the following:\\[1em]
JUDGMENT\\[0.5em]
in the civil case registered under No. 4567/2024\\
between\\[0.5em]
Tizio Rossi, residing in Turin, represented and defended by Attorney Lucia Verdi, with elected domicile at her office in Turin -- plaintiff\\[0.5em]
and\\[0.5em]
Caio Neri, residing in Milan, represented and defended by Attorney Marco Blu, with elected domicile at his office in Milan -- defendant\\[1em]
FACTS\\
The plaintiff summoned the defendant to court, alleging that on January 10, 2023, he lent him the sum of EUR 15,000, as evidenced by a private written agreement signed by both parties, with the obligation to repay by June 30, 2023, pursuant to Article 1813 of the Italian Civil Code (loan contract). Despite repeated requests, the defendant allegedly failed to repay the sum. The defendant, appearing in court at the Court Building in Turin, contested the existence of the debt, claiming that the sum received was a gift under Article 769 of the Italian Civil Code.\\[1em]
REASONS FOR THE DECISION\\
The evidence confirmed the existence of a loan contract under Article 1813 of the Italian Civil Code, as established by the private written agreement submitted in evidence and by the testimony of two independent witnesses (pursuant to Articles 2721--2725 of the Italian Civil Code on testimonial evidence). The claim of a gift is unfounded, as no donative intent on the part of the plaintiff was proven, as required by Article 769 of the Italian Civil Code. Pursuant to Article 2697 of the Italian Civil Code, the burden of proof lies with the party asserting release from an obligation, and in this case the defendant failed to provide adequate proof. Therefore, the plaintiff's credit must be recognized in the amount of EUR 15,000, plus statutory interest from July 1, 2023, in accordance with Article 1282 of the Italian Civil Code.\\[1em]
RULING\\
The Court, definitively ruling, dismissing any other claims or objections,\\
orders Mr. Caio Neri to pay Mr. Tizio Rossi the sum of EUR 15,000, plus statutory interest from July 1, 2023 until payment in full (Articles 1282 and 1284 of the Italian Civil Code);\\
orders the defendant to reimburse the plaintiff's legal costs, set at EUR 2,500 for legal fees, plus statutory surcharges, pursuant to Articles 91 and 92 of the Italian Code of Civil Procedure.\\[1em]
Turin, March 10, 2025\\
The Judge\\
Andrea Bianchi
\end{minipage}
\caption{Example of an abbreviated Italian civil judgment with references to legal articles and location entities.}
\label{fig:civil-judgment-legal-articles}
\end{figure}

